% 说明表与相同内容普通段落的纵向边界；不以 Word 数学字形为基准。
%! branches: stage=final,degree-level=doctor|bachelor,campus=harbin,lang=zh
%! engines: xelatex,lualatex
%! options: fontset=fandol
\newcommand\EqBoundaryRows[1]{%
  \begin{eqdenote}
  \item x & #1One 中文第一行；
  \item x & #1Two 中文第二行；
  \item x & #1Three 中文第三行。
  \end{eqdenote}}
\newcommand\EqBoundaryPlain[1]{%
  \noindent 式中 x——#1One 中文第一行；\par
  \noindent x——#1Two 中文第二行；\par
  \noindent x——#1Three 中文第三行。\par}
\newcommand\EqBoundaryEnd[1]{%
  \noindent #1After 中文接排 Agjp\par
  \section{#1Heading 中文标题}[#1 heading]
  \noindent #1End 中文结束 Agjp\par}
\newcommand\EqBoundaryGap{\par\vspace{20bp}}
\begin{document}
\mainmatter
\noindent BodyTableBefore 中文边界 Agjp\par
\EqBoundaryRows{BodyTable}\EqBoundaryEnd{BodyTable}
\EqBoundaryGap
\noindent BodyPlainBefore 中文边界 Agjp\par
\EqBoundaryPlain{BodyPlain}\EqBoundaryEnd{BodyPlain}
\clearpage
\noindent OpenTableBefore 中文边界 Agjp%
\EqBoundaryRows{OpenTable}\EqBoundaryEnd{OpenTable}
\EqBoundaryGap
\noindent OpenPlainBefore 中文边界 Agjp\par
\EqBoundaryPlain{OpenPlain}\EqBoundaryEnd{OpenPlain}
\clearpage
\begin{equation}x=1\quad\text{SmallTableBefore}\end{equation}
\EqBoundaryRows{SmallTable}\EqBoundaryEnd{SmallTable}
\EqBoundaryGap
\begin{equation}x=1\quad\text{SmallPlainBefore}\end{equation}
\EqBoundaryPlain{SmallPlain}\EqBoundaryEnd{SmallPlain}
\clearpage
\begin{equation}x=\frac{1+\psi^3}{1-\psi}\quad\text{LargeTableBefore}\end{equation}
\EqBoundaryRows{LargeTable}\EqBoundaryEnd{LargeTable}
\EqBoundaryGap
\begin{equation}x=\frac{1+\psi^3}{1-\psi}\quad\text{LargePlainBefore}\end{equation}
\EqBoundaryPlain{LargePlain}\EqBoundaryEnd{LargePlain}
\clearpage
\section{HeadTableBefore 中文标题}[Table heading boundary]
\EqBoundaryRows{HeadTable}\EqBoundaryEnd{HeadTable}
\EqBoundaryGap
\section{HeadPlainBefore 中文标题}[Paragraph heading boundary]
\EqBoundaryPlain{HeadPlain}\EqBoundaryEnd{HeadPlain}
\clearpage
\noindent ChainTableBefore 中文边界 Agjp\par
\begin{eqdenote}\item x & ChainTableOne 中文第一行；\end{eqdenote}
\begin{eqdenote}[]\item x & ChainTableTwo 中文第二行；\end{eqdenote}
\begin{eqdenote}[]\item x & ChainTableThree 中文第三行。\end{eqdenote}
\EqBoundaryEnd{ChainTable}
\EqBoundaryGap
\noindent ChainPlainBefore 中文边界 Agjp\par
\EqBoundaryPlain{ChainPlain}\EqBoundaryEnd{ChainPlain}
\clearpage
\EqBoundaryRows{TopTable}\EqBoundaryEnd{TopTable}
\clearpage
\EqBoundaryPlain{TopPlain}\EqBoundaryEnd{TopPlain}
\end{document}
